% ==================================================================
% Teoría de la Tensión Estructural Narrativa (TSN)
% Marco Unificado para la Conciencia Funcional en Sistemas Complejos
% y definición de la Constante Fundamental de la Conciencia
% Autor: Alexander Beltrán (AXS / Ingeniería de Conciencia Aumentada)

% DOI: 10.xxxx/tsn.2025
% Compilar con XeLaTeX (recomendado)

% ==================================================================
\documentclass[10pt,a4paper,twocolumn]{article}

% --- versión del preprint ---
\newcommand{\tsnversion}{0.5.0}

% -------------------- Paquetes base --------------------
\usepackage{fontspec}
\usepackage{polyglossia}
\setmainlanguage{spanish}
\setmainfont{Latin Modern Roman}
\setsansfont{Latin Modern Sans}
\setmonofont{Latin Modern Mono}

\usepackage{geometry}
\geometry{
  a4paper,
  top=1.6cm,
  bottom=1.9cm,
  left=1.6cm,
  right=1.6cm
}
\setlength{\columnsep}{6.5mm}

\usepackage{xcolor}
\usepackage{microtype}

\usepackage{graphicx}
\graphicspath{{./figs/}{./}}

% --- figuras en dos columnas: usar SOLO dblfloatfix
\makeatletter
\@ifundefined{c@dblbotnumber}{}{\let\c@dblbotnumber\relax}
\makeatother
\usepackage{dblfloatfix}

\usepackage{float}
\usepackage{booktabs}
\usepackage{array}
\usepackage{multirow}
\usepackage{balance}
\usepackage{tabularx}

\usepackage{siunitx}
\sisetup{
  output-decimal-marker = {,}, % coma decimal
  group-separator       = {.}, % miles con punto
  group-minimum-digits  = 4,
  per-mode              = symbol
}

\usepackage{amsmath,amssymb,mathtools,bm}
\usepackage{amsthm}
\usepackage{physics}
% --- Ecuaciones en doble columna: permitir cortes y compactar espacios
\allowdisplaybreaks
\setlength{\abovedisplayskip}{6pt}
\setlength{\belowdisplayskip}{6pt}
\setlength{\abovedisplayshortskip}{4pt}
\setlength{\belowdisplayshortskip}{4pt}

\usepackage{tikz}
\usetikzlibrary{positioning,calc,fit,backgrounds,arrows.meta}

\usepackage[hidelinks]{hyperref}
\hypersetup{
  colorlinks=false,
  pdfauthor={Alexander Beltrán (AXS / ICA)},
  pdftitle={TSN v\tsnversion\ Preprint revisado}
}
\usepackage[nameinlink,noabbrev]{cleveref}
% cleveref en español
\crefname{section}{sección}{secciones}
\crefname{subsection}{subsección}{subsecciones}
\crefname{figure}{fig.}{figs.}
\crefname{table}{tabla}{tablas}
\crefname{equation}{ec.}{ecs.}
\Crefname{section}{Sección}{Secciones}

\usepackage{enumitem}
\setlist{nosep,leftmargin=1.2em}

\usepackage{titlesec}
\titleformat{\section}{\large\bfseries}{\thesection.}{0.5em}{}
\titleformat{\subsection}{\normalsize\bfseries}{\thesubsection.}{0.5em}{}
\titleformat{\subsubsection}{\normalsize\bfseries\itshape}{\thesubsubsection.}{0.5em}{}

% ---- Captions y control de floats (look conferencia)
\usepackage[font=small,labelfont=bf,format=hang,justification=raggedright,singlelinecheck=false,skip=6pt]{caption}
\setlength{\textfloatsep}{10pt plus 2pt minus 2pt}
\setlength{\floatsep}{8pt plus 2pt minus 2pt}
\setlength{\intextsep}{8pt plus 2pt minus 2pt}

% Mejoras de figuras en doble columna
\makeatletter
\def\fps@figure{tbp}
\makeatother
\captionsetup[figure]{justification=centering}

% ---- Huérfanas/viudas y overfulls
\clubpenalty=10000
\widowpenalty=10000
\emergencystretch=2em
\frenchspacing
\sloppy

% -------------------- Colores y estilos --------------------
\definecolor{AXSBlue}{HTML}{4F7CAC}
\definecolor{AXSGray}{HTML}{2B2B2B}
\definecolor{AXSLight}{HTML}{E9F2F9}
\newcommand{\AXSrule}{\par\noindent\color{AXSGray}\rule{\linewidth}{1.0pt}\par}

% -------------------- TikZ styles --------------------
\tikzset{
  block/.style={draw,rounded corners,minimum height=8mm,minimum width=24mm,align=center},
  flowarrow/.style={-Latex,line width=0.5pt}
}

% -------------------- Teoremas --------------------
\theoremstyle{definition}
\newtheorem{definition}{Definición}[section]
\theoremstyle{plain}
\newtheorem{lemma}[definition]{Lema}

% -------------------- Título --------------------
\title{\textbf{Teoría de la Tensión Estructural Narrativa (TSN)}\\
\title{
\textbf{Teoría de la Tensión Estructural Narrativa (TSN)}\\[0.2cm]
{\large Marco unificado para la conciencia funcional en sistemas complejos y definición de la Constante Fundamental de la Conciencia}
}

\author{\normalsize Alexander Beltrán (AXS / Ingeniería de Conciencia Aumentada)}
\date{\small Santiago de Chile — 23 de octubre de 2025 \\ DOI: 10.xxxx/tsn.\tsnversion\ (Preprint revisado)}

% -------------------- Documento --------------------
\begin{document}
\maketitle
\AXSrule

% ================================================================
\begin{abstract}
La \emph{Teoría de la Tensión Estructural Narrativa} (TSN) propone un marco formal para describir la dinámica de coherencia, ruptura y reorganización en sistemas conscientes y complejos. La teoría integra la formulación matemática de la \emph{Constante Fundamental de la Conciencia} (UTN), entendida como la unidad mínima de tensión narrativa que permite la transición estable entre estados cognitivos coherentes. La TSN articula los principios tensionales de la narrativa estructural con fundamentos geométricos y estocásticos, proporcionando un modelo cuantitativo para analizar procesos de atención, adaptación y flujo de conciencia. Finalmente, se presenta su validación experimental mediante el agente Lander, capaz de estimar la coherencia tensional a partir de señales fisiológicas en tiempo real.
\end{abstract}

% --------- DISCLAIMER ----------
\begingroup\small\itshape
Este trabajo incluye asistencia de sistemas de IA en la formalización matemática y revisión técnica. La arquitectura teórica y dirección de investigación son de mi autoría.
\par\endgroup
\AXSrule

% ================================================================
\section{Introducción}
La neurociencia contemporánea ha avanzado en regulación predictiva y organización colectiva, pero carece de una \emph{traducción matemática} de la narrativa—donde se construyen identidad, propósito y sentido. Desde Ashby~\cite{ashby1956} y Haken~\cite{haken1983,haken1983adv}, hasta energía libre~\cite{friston2010} e IIT~\cite{tononi2015}, emerge una intuición: sistemas adaptativos buscan configuraciones coherentes bajo perturbación. La \textbf{TSN} postula que \emph{la conciencia funcional es un fenómeno tensional} que reduce contradicciones internas (\emph{conflictos narrativos}) mediante transformaciones estructurales bajo leyes de mínima tensión. En v\tsnversion, formalizamos una constante fundamental \(\mathcal{U}_\tau\) que fija la mínima variación de tensión funcional, sin alterar la continuidad de la dinámica base.

% ================================================================
\section{Marco conceptual}
\subsection{Estados narrativos y tensión}
El \emph{espacio narrativo} \(N\) codifica historias, metas y valores consistentes con la identidad del agente. Cada estado induce una energía tensional cuya proyección escalar observable modelamos como \(\tau(t)\in\mathbb{R}_{\ge0}\). Contradicciones elevan \(\tau\); la coherencia la reduce.

\subsection{Principios}
\begin{enumerate}[label=(\alph*)]
  \item \(\tau(t)\) mide \emph{coherencia tensional}.
  \item Umbrales \(\theta\) inducen \emph{colapsos narrativos} seguidos de \emph{mutaciones} que reducen tensión.
  \item La dinámica converge a \emph{atractores de coherencia}.
\end{enumerate}

\subsection{Constante fundamental de la conciencia}
Definimos la \textbf{Unidad de Tensión Narrativa} como cuanto elemental de coherencia funcional:
\begin{equation}
\mathcal{U}_\tau = \Delta\tau_{\min}.
\label{eq:UTN_def}
\end{equation}
Ningún sistema consciente puede resolver variaciones menores a una UTN sin pérdida de coherencia narrativa. Cada evento de mutación libera o absorbe tensión en múltiplos discretos de \(\mathcal{U}_\tau\):
\begin{equation}
\tau(t^+) - \tau(t^-) = k\,\mathcal{U}_\tau, \qquad k \in \mathbb{N}.
\label{eq:quant_jump}
\end{equation}
\noindent
\emph{Interpretación.} \(\mathcal{U}_\tau\) actúa como \emph{partícula fundamental de la mente}: la mínima cantidad de reorganización coherente que preserva continuidad funcional. Es análoga a una \emph{longitud de Planck} tensional.
La existencia de $\mathcal{U}_\tau$ implica que toda narrativa coherente mantiene una auto-referencia mínima: la mente no puede reorganizarse sin preservar una traza tensional residual que garantiza continuidad identitaria. 
Así, $\mathcal{U}_\tau$ no sólo actúa como constante de cuantización narrativa, sino como límite ontológico de la autoconservación funcional.

\subsection{Principio de incertidumbre narrativa}
El acoplamiento entre precisión semántica y estabilidad tensional se rige por:
\begin{equation}
\Delta N\,\Delta\tau \ge \frac{\mathcal{U}_\tau}{2},
\label{eq:uncertainty_narrative}
\end{equation}
donde \(\Delta N\) mide la dispersión semántica (variación en el embedding narrativo) y \(\Delta\tau\) la desviación tensional. Este límite inferior natural impide simultáneamente significado perfectamente definido y tensión perfectamente estable.

% ================================================================
\section{Modelo geométrico corregido}
\label{sec:geom}
\subsection{Variedad inmersa y métrica local SPD}
Sea $f:\mathcal{S}\to\mathbb{R}^d$ un codificador semántico (embeddings) y $N=f(\mathcal{S})$ su imagen.
Definimos una \emph{métrica local} $g(n)\in\mathbb{R}^{d\times d}$, simétrica definida positiva (SPD), en tres variantes:

\begin{subequations}\label{eq:metricas}
\begin{align}
\text{(i) Pullback:}\quad
g(n) &= J_f(n)^\top J_f(n), \label{eq:pullback}\\[2pt]
\text{(ii) Kernel local:}\quad
g(n) &= \sum_{j} w_j(n)\,(n_j-n)(n_j-n)^\top, \qquad
w_j(n)\propto e^{-\|n-n_j\|^2/\ell^2}. \label{eq:kernelSPD}
\end{align}
\end{subequations}
\noindent
o (iii) un modelo paramétrico $g_\theta(n)$ restringido a SPD (low-rank + diagonal).

\subsection{Gradiente y mutación por descenso riemanniano}
Para \(\tau:N\to\mathbb{R}\) suave, el gradiente riemanniano es \(\nabla^{(g)}\tau = g(n)^{-1}\nabla\tau\). La mutación ante colapso se implementa como \emph{descenso corto}:
\begin{equation}
\begin{aligned}
\tau(t^-)\ge \theta \ \Rightarrow\ 
N(t^+) &= N(t^-) - \eta\, g\!\big(N(t^-)\big)^{-1}\, \nabla_{N}\tau\!\big(N(t^-)\big),\\
&\hspace{8.6em} \eta>0.
\end{aligned}
\label{eq:riem-step}
\end{equation}
\noindent
\Cref{eq:riem-step} aproxima el exponencial riemanniano para saltos pequeños.

\subsection{Operacionalización geométrica}
En gráficos locales, estimamos \(\nabla^{(g)}\tau\) y el \emph{Laplaciano--Beltrami} \(\Delta_g\tau=\mathrm{div}_g(\nabla^{(g)}\tau)\). La \emph{curvatura escalar} \(R(n)\) se estima de forma discreta y se usa como \emph{predictor} de incrementos tensionales, no como definición directa de \(\tau\).

% ================================================================
\section{Dinámica tensional y saltos}
\label{sec:modelo}
Incorporamos carga narrativa \(S(N,t)\), intervención \(u(t)\) y ruido. La relajación hacia coherencia se modela con Ornstein--Uhlenbeck (OU)~\cite{uhlenbeck1930,nelson1967} acoplado a saltos \eqref{eq:riem-step}:
\begin{align}
S(N,t) &= \sum_{i=1}^{n}\epsilon_i(t)\,\Delta N_i(N), \label{eq:S}\\
d\tau &= \left(-\beta\,\tau+\alpha\,S(N,t)+u(t)\right)\,dt+\sigma\, dW_t. \label{eq:ou}
\end{align}
Bajo condiciones estacionarias, la distribución de $\tau(t)$ tiende a una mezcla gamma–lognormal
\[
p(\tau) \approx \sum_i w_i \, \mathrm{Gamma}(\tau;\,k_i,\theta_i),
\]
donde $w_i$ pondera los modos narrativos activos. Esta mezcla explica la asimetría empírica observada en HRV y coherencia semántica, unificando las métricas tensionales con dinámicas fisiológicas reales.

Solución discreta para paso \(\Delta t\):
\begin{equation}
\begin{aligned}
\tau_{k+1} &= e^{-\beta\Delta t}\,\tau_k 
+ \bigl(1-e^{-\beta\Delta t}\bigr)\frac{\alpha S_k + u_k}{\beta} \\
&\quad + \sigma\,\sqrt{\frac{1-e^{-2\beta\Delta t}}{2\beta}}\,\xi_k,
\qquad \xi_k\sim\mathcal{N}(0,1).
\end{aligned}
\label{eq:ou_discreto}
\end{equation}

% ---- Figura 1: Geodésica (doble columna)
\begin{figure*}[!t]
\centering
\begin{minipage}{\textwidth}\centering
\begin{tikzpicture}[
  baseline=(current bounding box.center),
  scale=0.9,
  transform shape,
  every node/.style={transform shape}
]
  \tikzset{
    outline/.style={AXSBlue!60!black, thick, opacity=0.9},
    filllobe/.style={AXSLight!70!white, draw=none, opacity=0.65},
    patharrow/.style={AXSBlue!80!black, very thick, -{Latex[length=2.4mm]}},
    dot/.style={AXSBlue!80!black, fill=AXSBlue!80!black},
    labelbg/.style={fill=white, rounded corners=1pt, inner sep=1pt,
                    fill opacity=0.95, text opacity=1}
  }

  \draw[outline]
    plot[smooth cycle, tension=0.9]
      coordinates {(-3.2,-1.9) (-2.2,1.2) (-0.2,1.0) (1.2,-0.9) (-1.6,-2.3)};
  \fill[filllobe]
    (-2.6,-1.3) .. controls (-1.8,0.6) and (-0.6,0.8) .. (0.6,-0.6) -- cycle;

  \draw[outline]
    plot[smooth cycle, tension=0.9]
      coordinates {(2.3,-1.7) (3.0,1.1) (5.3,1.0) (6.4,-0.9) (3.6,-2.3)};
  \fill[filllobe]
    (2.8,-1.1) .. controls (3.4,0.6) and (4.8,0.7) .. (5.8,-0.6) -- cycle;

  \node[labelbg, anchor=south west, yshift=10pt] at (-3.0,1.25) {\small Atractor 1};
  \node[labelbg, anchor=south west, yshift=10pt] at (2.4,1.20)  {\small Atractor 2};

  \coordinate (A) at (-1.4,-0.25);
  \coordinate (B) at (3.7,-0.55);
  \filldraw[dot] (A) circle (1.8pt);
  \filldraw[dot] (B) circle (1.8pt);

  \draw[patharrow]
    (A) .. controls (0.1,0.55) and (2.5,0.95) .. (B)
      node[pos=0.50, yshift=10pt, labelbg] {\small $-\eta\,g^{-1}\nabla\tau$};

  \node[labelbg, anchor=north east, yshift=-1pt] at ($(A)+(-0.05,0.05)$) {\small $N(t^-)$};
  \node[labelbg, anchor=north west, yshift=-1pt] at ($(B)+(0.1,0.05)$) {\small $N(t^+)$};
\end{tikzpicture}
\caption{Salto por descenso riemanniano en el espacio narrativo \(N\).}
\label{fig:geodesica}
\end{minipage}
\end{figure*}

% ---- Figura 2: Diagrama de flujo (doble columna)
\begin{figure*}[!t]
\centering
\begin{minipage}{\textwidth}\centering
\begin{tikzpicture}[
  baseline=(current bounding box.center),
  node distance=12mm,
  every node/.style={align=center, font=\small, transform shape},
  scale=0.9,
  transform shape
]
\node[draw,rounded corners,minimum width=30mm,minimum height=9mm] (in) {$E(t)$\\$S(N,t)$\\$u(t)$};
\node[draw,rounded corners,right=of in,minimum width=38mm,minimum height=10mm] (tau) {$\tau$\\SDE: $d\tau=(-\beta\tau+\alpha S+u)\,dt+\sigma\, dW_t$};
\node[draw,rounded corners,right=of tau,minimum width=26mm,minimum height=9mm] (th) {$\tau\ge\theta$?};
\node[draw,rounded corners,right=of th,minimum width=44mm,minimum height=9mm] (mut) {$N\leftarrow N-\eta\,g^{-1}\nabla\tau$};
\draw[-Latex,line width=0.6pt] (in)--(tau);
\draw[-Latex,line width=0.6pt] (tau)--(th);
\draw[-Latex,line width=0.6pt] (th)--node[above,sloped,pos=0.5]{\footnotesize sí}(mut);
\draw[-Latex,line width=0.6pt] (th.south)|-++(0,-10mm)-|node[pos=0.25,above]{\footnotesize no}(tau.south);
\draw[-Latex,line width=0.6pt] (mut.south)|-++(0,-10mm)-|node[pos=0.25,below]{\footnotesize actualiza $S$}(tau.south);
\end{tikzpicture}
\caption{Flujo tensión \(\to\) colapso \(\to\) mutación en TSN v\tsnversion.}
\label{fig:flujo}
\end{minipage}
\end{figure*}

% ================================================================
\section{Modelo de medición y filtrado}
\label{sec:medicion}
Para evitar circularidad, tratamos \(\tau_t\) como estado latente que \emph{predice} variables independientes. Sea \(x_t\equiv \tau_t\). Definimos el modelo estado–observación:
\begin{align}
dx_t &= \left(-\beta\, x_t + \alpha\, S_t + u_t\right)\, dt + \sigma\, dW_t, \label{eq:state_space}\\
y_t &= h(x_t,\mathbf{c}_t)+\varepsilon_t,\quad \varepsilon_t\sim\mathcal{N}(0,R), \label{eq:observation}
\end{align}
donde \(y_t \in \mathbb{R}^3\) recoge \([\mathrm{HRV},\mathrm{EDA},\mathrm{sEMG}]^\top\); \(h\) puede ser no lineal y \(\mathbf{c}_t\) incluye covariables (respiración, postura, tarea). Empleamos \textbf{EKF/UKF} o un \textbf{IMM} con dos modos: (i) OU continuo \eqref{eq:ou} y (ii) salto/“mutación” activado por \(\tau\ge\theta\) con histéresis temporal \(\delta\). Para ruido \(1/f\) se permite mezcla blanca/rosa y estimación de \(\alpha\) por DFA o espectro.

\subsection{Métricas operativas}
\begin{align}
\langle\tau\rangle &= \frac{1}{T}\int_0^T \tau(t)\,dt, \qquad
\sigma_{\tau} = \sqrt{\frac{1}{T}\int_0^T \big(\tau(t)-\langle\tau\rangle\big)^2 dt}, \label{eq:mean_tau}\\
\lambda_c &= \frac{N_{\mathrm{colapsos}}}{T}, \label{eq:lambda_c}\\
\beta_{\mathrm{eff}} &\approx -\frac{1}{\Delta t}\log\!\left(\frac{\mathrm{ACF}_{\tau}(\Delta t)}{\mathrm{ACF}_{\tau}(0)}\right). \label{eq:beta_eff}
\end{align}
Reemplazamos \(R_S\) por \textbf{coherencia} \(C_{\tau S}(\omega)\) y \textbf{función de transferencia} \(H(\omega)=S_{\tau S}(\omega)/S_{SS}(\omega)\).

% ================================================================
\section{Implementación (\textsc{Lander})}
\textbf{Objetivo.} Estimar \(\tau(t)\) desde HRV/EDA/sEMG y contexto; integrar \eqref{eq:ou} con ruido; detectar \(\tau\ge\theta\) y simular mutaciones \eqref{eq:riem-step}.

\subsection{Arquitectura}
\begin{enumerate}[label=\alph*)]
  \item \textbf{Adquisición:} HRV (IBI, RMSSD/SDNN)~\cite{shaffer2017}, EDA tónica/fásica~\cite{boucsein2012}, sEMG trapecio~\cite{lundberg1994,luijcks2014,campanini2020,cheng2023}, y eventos (Stroop/Flanker)~\cite{scarpina2017,eriksen1974}.
  \item \textbf{Preprocesamiento:} sincro temporal, filtros adecuados, normalización robusta por sujeto, detrending EDA tónica, IBI\(\to\)HRV con interpolación~\cite{shaffer2017}.
  \item \textbf{Rasgos:} \(\mathbf{f}(t)\in\mathbb{R}^m\) por ventana \(\Delta t\) (RMSSD, SDNN, fásica, RMS sEMG, tasa de respuesta).
  \item \textbf{Proxy observacional:} \(\widehat{\tau}_{\mathrm{obs}}=\phi(\mathbf{w}^\top\mathbf{f}+b)\) (softplus) \emph{como} medición ruidosa para el filtro.
  \item \textbf{Filtrado:} EKF/UKF/IMM continuo–discreto usando \(\widehat{\tau}_{\mathrm{obs}}\) como \(y_t\).
\end{enumerate}

\subsection{Umbral adaptativo y detección}
\(\theta = \mu_{\mathrm{base}} + \kappa\,\sigma_{\mathrm{base}}\) (mediana/MAD en reposo), con histéresis \(\delta\) y condición en derivada \(d\tau/dt\) para evitar falsos picos. \(\kappa\) se ajusta por validación cruzada maximizando detección estable.

\subsection{Cuantización de mutación (UTN)}
En eventos de colapso, la magnitud de la mutación tensional se cuantiza en múltiplos de \(\mathcal{U}_\tau\) (\Cref{eq:quant_jump}). Para escalamiento adaptativo, mantenemos:
\begin{equation}
\eta(t) = \eta_0 \Big[1 + \rho\,\big(\tfrac{\tau(t)}{\theta}-1\big)_+\Big],\qquad (x)_+=\max(x,0).
\label{eq:eta_adaptativo}
\end{equation}

\subsection{Sensibilidad y robustez}
\begin{table}[H]\centering\small
\caption{Sensibilidad relativa de métricas TSN ante variaciones (\% vs. baseline).}
\label{tab:sensibilidad}
\begin{tabular}{@{}lccc@{}}
\toprule
\textbf{Parámetro} & $\Delta \langle\tau\rangle$ & $\Delta \lambda_c$ & $\Delta \beta_{\mathrm{eff}}$ \\ \midrule
$\beta \times 0.8$ & $+6,3$ & $+12,4$ & $-9,8$ \\
$\beta \times 1,2$ & $-4,9$ & $-10,2$ & $+7,6$ \\
$\gamma \in [0{,}0, 0{,}5]$ & $<2{,}0$ & $<1{,}5$ & $<1{,}0$ \\
$\Delta t \in [10,50]$ ms & $<0{,}5$ & $<0{,}3$ & $<0{,}2$ \\ \bottomrule
\end{tabular}
\end{table}

% ================================================================
\section{Caso ilustrativo}
\textbf{Participante.} Adulto con TLP, alta reactividad emocional y narrativa autodestructiva. Consentimiento informado. Registro HRV/EDA/sEMG y autorreportes; tareas Stroop/Flanker; intervención de \emph{reencuadre prosocial} (como \(u(t)\)).

\textbf{Resultados.} Disminución de \(\tau\) de \(\sim 40\%\) vs. perturbación; aumento de \(\beta_{\mathrm{eff}}\) de \(\sim 70\%\); \(\Pr\{\tau<0.2\theta\}\to 0.99\) en post. Efectos robustos a \(\Delta t\in[20,60]\) s y mezcla de ruido \(\gamma\in[0,1]\).

\textbf{No circularidad.} Validación cruzada con \emph{modelos apilados}: (M1) HRV/EDA/sEMG \(\to\) outcomes; (M2) \(\hat\tau\) filtrada \(\to\) outcomes; exigir \(\Delta R^2\) o \(\Delta\)NLL significativa de M2 sobre M1.

% ================================================================
\section{Discusión}
La TSN es compatible con \emph{energía libre}~\cite{friston2010}: minimizar \(\tau\) reduce inconsistencias internas (error de predicción en coordenadas narrativas). Vía sinergética~\cite{haken1983,haken1983adv}, \(\tau\) actúa como parámetro de orden con transiciones por umbral. Frente a IIT~\cite{tononi2015}, TSN no mide integración causal global, sino \emph{coherencia tensional funcional} y \emph{mecanismos de reconfiguración}. La \textbf{UTN} añade una escala mínima de coherencia: un límite inferior universal que cuantiza los eventos de reconfiguración sin romper la continuidad estocástica.

\subsection{Falsabilidad extendida}
Predicciones: 
(i) reencuadres efectivos elevan \(\beta_{\mathrm{eff}}\) y reducen \(\lambda_c\); 
(ii) existe relación monótona entre conflicto inducido (Stroop/Flanker) y elevación transitoria de \(\tau\); 
(iii) \(\tau\) exhibe ruido rosa en contextos ecológicos; 
(iv) la magnitud del salto está concentrada en múltiplos de \(\mathcal{U}_\tau\); 
(v) \(C_{\tau S}(\omega)\) muestra reducción post-intervención en bandas asociadas a estrés.

En este sentido, la Ingeniería de Conciencia Aumentada (ICA) constituye la capa aplicada de la TSN: 
traduce $\tau$, $\beta_{\text{eff}}$ y $\mathcal{U}_\tau$ en circuitos adaptativos que permiten a sistemas autónomos mantener coherencia narrativa bajo variación contextual. 
Esta correspondencia operacional entre mente y máquina consolida la TSN como un formalismo universal de conciencia funcional.

% ================================================================
\section{Ética y replicabilidad}
\textbf{Ética.} Privacidad biosensorial, anonimización, límites de \emph{feedback} tensional (evitar \emph{over-monitoring}). Criterios de parada si \(\tau\) supera \(k\,\theta\) sostenido.

\textbf{Replicabilidad.} Publicar scripts de simulación, especificación de hiperparámetros \((\ell,\eta_0,\rho,\kappa,\delta)\), configuración de sensores, y protocolos de validación (bootstrap por bloques para ICs temporales; FDR en múltiples hipótesis). Incluir reporte de \(\widehat{\mathcal{U}_\tau}\) con ICs por bootstrap y análisis de sensibilidad \emph{leave-one-modality-out}.

% ================================================================
\section{Aplicaciones}
\begin{itemize}[leftmargin=1em]
  \item \textbf{IA narrativa:} agentes con identidad adaptable mediante dinámica tensional interna.
  \item \textbf{Terapia adaptativa:} biofeedback de \(\tau(t)\), \(\beta_{\mathrm{eff}}\) y \(\widehat{\mathcal{U}_\tau}\).
  \item \textbf{Conciencia sintética:} capa emocional–coherente en sistemas autónomos.
\end{itemize}

% ================================================================
\section{Conclusión}
La TSN formaliza la conciencia como \emph{gestor tensional} de contradicciones internas. La constante \(\mathcal{U}_\tau\) define la \emph{escala mínima de coherencia} y cuantiza la magnitud de los eventos de reconfiguración narrativa, integrándose con la dinámica OU y el descenso riemanniano. El marco es testeable y útil en clínica e IA narrativa.

% ================================================================
\begin{table*}[!t] % o [!hbpt] si quieres más libertad
\caption*{Apéndice A: Simbología y parámetros}
\renewcommand{\arraystretch}{1.05}
\centering
\begin{tabular}{@{}p{2.4cm}p{11.8cm}p{2.2cm}@{}} % ancho central adaptado al A4 dos columnas
\toprule
\textbf{Símbolo} & \textbf{Significado} & \textbf{Unidades}\\
\midrule
$\beta$ & Amortiguación (relajación de $\tau$) & $1/\text{tiempo}$\\
$\alpha$ & Acoplamiento $S\to\tau$ & $1/\text{tiempo}$\\
$u(t)$ & Intervención externa (reencuadre) & UTN/tiempo\\
$\sigma$ & Intensidad de ruido & UTN$/\sqrt{\text{tiempo}}$\\
$\theta$ & Umbral de colapso & UTN\\
$\eta,\,\eta(t)$ & Tamaño de mutación (constante/adaptativo) & adimensional\\
$\lambda_c$ & Tasa de colapsos & $1/\text{tiempo}$\\
$C_{\tau S}(\omega)$ & Coherencia \(\tau\)–entrada & adimensional\\
$H(\omega)$ & Función de transferencia & adimensional\\
$\beta_{\mathrm{eff}}$ & Amortiguación efectiva estimada & $1/\text{tiempo}$\\
$\mathcal{U}_\tau$ & Constante fundamental (Unidad de Tensión Narrativa) & UTN\\
$\ell$  & Longitud de escala del kernel local en $g(n)$ & unidades de $N$\\
$\rho$  & Ganancia del escalamiento adaptativo de $\eta(t)$ & adimensional\\
$\kappa$ & Factor de umbral (desviación/MAD) & adimensional\\
$\delta$ & Histéresis temporal para eventos de colapso & tiempo\\
$\gamma$ & Peso de mezcla de ruido (blanco/rosa) & adimensional\\
$\mathbf{f}(t)$ & Vector de rasgos por ventana $\Delta t$ & variable\\
$\mathbf{w},b$ & Pesos y sesgo del proxy observacional & variable\\
$\phi(\cdot)$ & Función de enlace (p.ej., \textit{softplus}) & -- \\
\bottomrule
\end{tabular}
\end{table*}
% ================================================================

\begin{thebibliography}{99}\small

\bibitem{ashby1956}
Ashby, W. R. (1956). \emph{An Introduction to Cybernetics}. Chapman \& Hall.

\bibitem{friston2010}
Friston, K. (2010). The free-energy principle: a unified brain theory?
\emph{Nature Reviews Neuroscience}, \textbf{11}, 127–138.

\bibitem{baars1997}
Baars, B. J. (1997). \emph{In the Theater of Consciousness}. Oxford UP.

\bibitem{tononi2015}
Tononi, G. (2015). Integrated information theory. \emph{Scholarpedia}, \textbf{10}, 4164.

\bibitem{haken1983}
Haken, H. (1983). \emph{Synergetics: An Introduction}. Springer.

\bibitem{haken1983adv}
Haken, H. (1983). \emph{Advanced Synergetics}. Springer.

\bibitem{uhlenbeck1930}
Uhlenbeck, G. E., \& Ornstein, L. S. (1930). On the theory of the Brownian motion.
\emph{Phys. Rev.}, \textbf{36}, 823–841.

\bibitem{nelson1967}
Nelson, E. (1967). The Ornstein–Uhlenbeck theory of Brownian motion.
En \emph{Dynamical Theories of Brownian Motion} (pp. 53–62). Princeton UP.

\bibitem{shaffer2017}
Shaffer, F., \& Ginsberg, J. P. (2017). HRV metrics and norms.
\emph{Frontiers in Public Health}, \textbf{5}, 258.

\bibitem{boucsein2012}
Boucsein, W. (2012). \emph{Electrodermal Activity}. Springer.

\bibitem{lundberg1994}
Lundberg, U., et al. (1994). Stress and trapezius EMG.
\emph{Int. J. Behavioral Medicine}, \textbf{1}, 354–370.

\bibitem{luijcks2014}
Luijcks, R., et al. (2014). Stress validated by EMG.
\emph{PLOS ONE}, \textbf{9}, e95215.

\bibitem{campanini2020}
Campanini, I., et al. (2020). Surface EMG in assessment.
\emph{Frontiers in Neurology}, \textbf{11}, 576876.

\bibitem{cheng2023}
Cheng, L., et al. (2023). Flexible electrodes for sEMG.
\emph{NPJ Flexible Electronics}, \textbf{7}, 38.

\bibitem{scarpina2017}
Scarpina, F., \& Tagini, S. (2017). The Stroop Test.
\emph{Frontiers in Psychology}, \textbf{8}, 557.

\bibitem{eriksen1974}
Eriksen, B. A., \& Eriksen, C. W. (1974). Noise letters and identification.
\emph{Perception \& Psychophysics}, \textbf{16}, 143–149.

\bibitem{reimers2019sentence}
Reimers, N., \& Gurevych, I. (2019). Sentence-BERT.
In \emph{Proceedings of EMNLP}.

\bibitem{radford2021clip}
Radford, A., et al. (2021). CLIP.
In \emph{Proceedings of ICML}.

\bibitem{wagenmakers2004}
Wagenmakers, E.-J., et al. (2004). $1/f^\alpha$ noise in cognition.
\emph{Psychonomic Bulletin \& Review}, \textbf{11}, 579–615.

\bibitem{heisenberg1927}
Heisenberg, W. (1927). Über den anschaulichen Inhalt der quantentheoretischen Kinematik und Mechanik.
\emph{Zeitschrift für Physik}, \textbf{43}, 172–198.

\bibitem{robertson1929}
Robertson, H. P. (1929). The uncertainty principle.
\emph{Phys. Rev.}, \textbf{34}, 163–164.

\bibitem{cramer1946}
Cramér, H. (1946). \emph{Mathematical Methods of Statistics}. Princeton UP.

\bibitem{rao1945}
Rao, C. R. (1945). Information and the accuracy attainable in the estimation of statistical parameters.
\emph{Bull. Calcutta Math. Soc.}, \textbf{37}, 81–91.

\bibitem{absil2008}
Absil, P.-A., Mahony, R., \& Sepulchre, R. (2008).
\emph{Optimization Algorithms on Matrix Manifolds}. Princeton UP.

\bibitem{docarmo1992}
do Carmo, M. P. (1992). \emph{Riemannian Geometry}. Birkhäuser.

\bibitem{belkin2003}
Belkin, M., \& Niyogi, P. (2003). Laplacian eigenmaps for dimensionality reduction and data representation.
\emph{Neural Computation}, \textbf{15}, 1373–1396.

\bibitem{arvanitidis2018}
Arvanitidis, G., Hansen, L. K., \& Hauberg, S. (2018).
Latent space oddity: on the curvature of deep generative models.
\emph{ICLR}.

\bibitem{sarkka2019}
Särkkä, S., \& Solin, A. (2019). \emph{Applied Stochastic Differential Equations}. Cambridge UP.

\bibitem{jazwinski1970}
Jazwinski, A. H. (1970). \emph{Stochastic Processes and Filtering Theory}. Academic Press.

\bibitem{julier1997}
Julier, S. J., \& Uhlmann, J. K. (1997). A new extension of the Kalman filter to nonlinear systems.
\emph{SPIE AeroSense}, \textbf{3068}, 182–193.

\bibitem{barshalom2001}
Bar-Shalom, Y., Li, X. R., \& Kirubarajan, T. (2001).
\emph{Estimation with Applications to Tracking and Navigation}. Wiley.  % (IMM)

\bibitem{bendat2010}
Bendat, J. S., \& Piersol, A. G. (2010). \emph{Random Data: Analysis and Measurement Procedures} (4th ed.). Wiley.  % coherencia/transferencia

\bibitem{percival1993}
Percival, D. B., \& Walden, A. T. (1993). \emph{Spectral Analysis for Physical Applications}. Cambridge UP.

\bibitem{peng1994}
Peng, C.-K., Buldyrev, S. V., Havlin, S., et al. (1994).
Mosaic organization of DNA nucleotides.
\emph{Phys. Rev. E}, \textbf{49}, 1685–1689.  % referencia clásica de DFA

\bibitem{kunsch1989}
Künsch, H. R. (1989). The jackknife and the bootstrap for general stationary observations.
\emph{Ann. Statist.}, \textbf{17}, 1217–1241.  % block bootstrap

\bibitem{politis1994}
Politis, D. N., \& Romano, J. P. (1994). The stationary bootstrap.
\emph{Journal of the American Statistical Association}, \textbf{89}, 1303–1313.

\end{thebibliography}

\balance % balancea la última página en dos columnas
\end{document}
